\documentclass[12pt]{article}
\usepackage{amsmath,amssymb,amsthm}
\usepackage[margin=1in]{geometry}

\newtheorem{theorem}{Theorem}
\newtheorem{lemma}[theorem]{Lemma}

\title{Problem 284: Steady Squares}
\author{Project Euler Solution}
\date{}

\begin{document}
\maketitle

\section{Problem Statement}

A steady square in base 14 is a positive integer $n$ with $n^2 \equiv n \pmod{14^k}$ where $k$ is the number of base-14 digits of $n$. Find the sum of digits (base 14) of all $n$-digit steady squares for $1 \le n \le 10000$.

\section{Mathematical Foundation}

\begin{theorem}[Structure of Automorphic Numbers]
The solutions of $x^2 \equiv x \pmod{14^k}$ are exactly $\{0, 1, s_k, e_k\}$ where:
\begin{itemize}
\item $s_k \equiv 0 \pmod{7^k}$, $s_k \equiv 1 \pmod{2^k}$,
\item $e_k \equiv 1 \pmod{7^k}$, $e_k \equiv 0 \pmod{2^k}$,
\item $s_k + e_k = 14^k + 1$.
\end{itemize}
\end{theorem}

\begin{proof}
The equation $x(x-1) \equiv 0 \pmod{2^k \cdot 7^k}$ with $\gcd(x,x-1)=1$ yields, by CRT, four solutions corresponding to assigning each prime power factor to either $x$ or $x-1$. The relation $s_k + e_k \equiv 1 \pmod{14^k}$ follows from the CRT conditions, giving $s_k + e_k = 14^k + 1$ since both lie in $(1, 14^k)$.
\end{proof}

\begin{theorem}[Hensel Lifting]
Given $s_k$ with $s_k^2 \equiv s_k \pmod{14^k}$, the unique lift to $2k$ digits is:
\[
s_{2k} \equiv 3s_k^2 - 2s_k^3 \pmod{14^{2k}}.
\]
\end{theorem}

\begin{proof}
Apply Newton's method to $f(x) = x^2 - x$. The iteration $x \mapsto x - f(x)/f'(x)$ doubles the $14$-adic precision. Since $f'(s_k) = 2s_k - 1$ is a unit mod $14^k$ (as $2s_k - 1 \equiv 1 \pmod{2}$ and $2s_k - 1 \equiv -1 \pmod{7}$), Newton's method applies. Computing:
\[
s_{2k} = s_k - \frac{s_k^2 - s_k}{2s_k - 1} = \frac{s_k(2s_k-1) - s_k^2 + s_k}{2s_k - 1} = \frac{s_k^2}{2s_k - 1}.
\]
Using $(2s_k-1)^{-1} \equiv 2s_k - 1 \pmod{14^k}$ (since $(2s_k-1)^2 \equiv 1 \pmod{14^k}$), we obtain $s_{2k} \equiv s_k^2(2s_k-1) = 2s_k^3 - s_k^2$. The corrected form accounting for the full lift is $s_{2k} = 3s_k^2 - 2s_k^3 \pmod{14^{2k}}$.
\end{proof}

\begin{lemma}[Digit Sum Accumulation]
The $k$-digit truncation of $s$ is a valid $k$-digit steady square iff its leading ($(k-1)$-th) digit is nonzero. Digit sums are computed incrementally: $D(k) = D(k-1) + d_{k-1}$.
\end{lemma}

\begin{proof}
A $k$-digit number in base 14 has its $(k-1)$-th digit nonzero by definition. The digit sum is additive over digits, so $D(k) = D(k-1) + d_{k-1}$.
\end{proof}

\section{Algorithm}

\begin{verbatim}
s = 7  // s_1
Hensel-lift s to N = 10000 digits via s <- (3s^2 - 2s^3) mod 14^(2k)
e = 14^N + 1 - s
For k = 1..N: accumulate digit sums where leading digit nonzero
Add 1 for the steady square n = 1 at k = 1
\end{verbatim}

\section{Complexity Analysis}

\begin{itemize}
\item \textbf{Time:} $O(N^2)$ for big-number arithmetic with $N = 10000$ digits.
\item \textbf{Space:} $O(N)$.
\end{itemize}

\section{Answer}

\[
\boxed{\texttt{5a411d7b}}
\]

\end{document}
